\documentclass[11pt,a4paper]{article}
\usepackage{amsmath,amssymb,amsfonts}
\usepackage{amsbsy}
\usepackage{amsthm}
\usepackage{geometry}
\usepackage{enumitem}
\usepackage{hyperref}

\newtheorem{theorem}{Theorem}[section]
\newtheorem{lemma}[theorem]{Lemma}
\newtheorem{proposition}[theorem]{Proposition}
\newtheorem{corollary}[theorem]{Corollary}
\theoremstyle{definition}
\newtheorem{definition}[theorem]{Definition}
\theoremstyle{remark}
\newtheorem{remark}[theorem]{Remark}

\newcommand{\Jnorm}{J_{\rm norm}}
\newcommand{\Jbound}{\frac{3\pi^2}{8}}
\newcommand{\Mnorm}{M_{\rm norm}}

\title{Discrete Dual Spacetime: Bounded Torsion Scalar and Finite Discrete Rotor Image}

\author{https://github.com/hypernumbernet}
\date{\today}

\begin{document}

\maketitle

\begin{abstract}
The dual spacetime theory rejects the continuum hypothesis entirely, positing that each massive particle carries a pair of intrinsically compactified Minkowski spacetimes encoded in the 16-real-dimensional biquaternion algebra isomorphic to $\mathrm{Cl}(3,1)$. In this companion paper, we formalize the discretization of the rapidity parameters arising from the dual map $X \mapsto Xi$. Cyclic (elliptic) rotor exponentials are $2\pi$-periodic; hyperbolic (boost) exponentials are nowhere $2\pi$-periodic, so placing boosts on a $2\pi$-torus is a modelling choice rather than a consequence of the exponential map. Admissibility---non-negative embedded rapidities with anti-synchronisation $\alpha_a+\beta_a\leq\pi/2$ per axis---is an additional hypothesis, not a consequence of discretizing the torus. On that cone the normalized torsion scalar satisfies $|\Jnorm|\leq 1$, equivalently $|J|\leq 3\pi^2/8$. The admissible locus itself is a product of three triangular cones of side $\lfloor N/4\rfloor$, of cardinality $T(\lfloor N/4\rfloor)^3$ with $T(K)=(K+1)(K+2)/2$. The occupied values of $\Jnorm$ form a finite subset of the rational lattice of spacing $16/(3N^2)$: the first nonzero slot $\Delta=\pm 1$ is attained for every $N\geq 4$, so the ultraviolet floor $|\Jnorm|\geq 16/(3N^2)$ is sharp, while at finite $N$ the bound interval still contains empty slots (already at $N=8$, $|\Delta|\in\{10,11\}$). What is finite is the image of $(\mathbb{Z}/N\mathbb{Z})^6$ under the torsion-rotor map, not the unit group of an unspecified integer biquaternion order. Same-axis generators commute, so a one-axis exponential factorises; distinct-axis generators do not, and the discrete rotor image is therefore not a Cartesian product of three one-axis images. The resulting Diophantine constraints on gravitational solutions provide a number-theoretic foundation for quantum gravity effects while preserving exact equivalence to General Relativity in the continuum limit.
\end{abstract}

\section{Introduction}
\label{sec:intro}

The dual spacetime theory, as presented in Ref.~\cite{main-paper}, achieves an exact algebraic reformulation of General Relativity as the teleparallel equivalent (TEGR) using biquaternions, while interpreting gravity and inertia as torsional mismatch between particle-intrinsic usual and dual spacetimes. A central philosophical tenet is the complete rejection of the spacetime continuum hypothesis: no shared differentiable manifold exists; spacetime degrees of freedom are demoted to internal attributes of particles.

This rejection naturally invites a discrete structure at the fundamental scale. The biquaternion algebra $\mathbb{H} \oplus \mathbb{H} \cong \mathrm{Cl}(3,1)$ already encodes compactness through the pseudoscalar $i$ ($i^2 = -1$) and the dual map $X \mapsto Xi$, which reverses the temporal arrow while preserving the Minkowski norm. Only the cyclic sector of the rotor exponential is intrinsically periodic; the hyperbolic sector is not.

In this paper, we develop this discretization. By placing the six rapidity coordinates on the finite torus $(\mathbb{Z}/N\mathbb{Z})^6$, where $N$ is a large integer related to the particle's compactification scale, we obtain:
\begin{itemize}
  \item A bounded normalized torsion scalar $|\Jnorm| \leq 1$ on admissible configurations (equivalently $|J| \leq 3\pi^2/8$), due to the sign asymmetry of the Killing form together with the extra admissible-cone hypotheses (non-negative rapidities and anti-synchronisation). Discretization of the torus alone does not imply the bound.
  \item A sharp discrete refinement $|\Jnorm|\leq (4\lfloor N/4\rfloor/N)^2$, with equality $\pm 1$ attained if and only if $4\mid N$. When $N<4$ the admissible cone is a singleton and $\Jnorm=0$.
  \item An explicit count of admissible points, $T(\lfloor N/4\rfloor)^3$, and a rational spectrum $\Jnorm=16\Delta/(3N^2)$ for an integer mismatch $\Delta$. For $N\geq 4$ the value $\Delta=\pm 1$ is attained, so the smallest nonzero $|\Jnorm|$ is exactly $16/(3N^2)$.
  \item Spectral holes inside the bound: the image of $\Delta$ is the three-fold sumset of triangle differences of squares, a proper subset of $\{-3K^2,\ldots,3K^2\}$ with $K=\lfloor N/4\rfloor$. At $N=8$ the slots $|\Delta|\in\{10,11\}$ are empty.
  \item A finite \emph{discrete rotor image}: the image of $(\mathbb{Z}/N\mathbb{Z})^6$ under $t\mapsto\exp(\Omega_{\mathrm{torsion}}(t))$. This is not the unit group of an integer biquaternion order, which may be infinite (Pell-type), and it is not the Cartesian product of three one-axis images, because distinct-axis generators do not commute.
  \item Diophantine equations governing allowable torsional configurations. Quantizing hyperbolic rapidities onto a $2\pi$-torus is a modelling choice; balanced-type seeds have vanishing modular winding, so a torus-winding obstruction cannot fire on that locus.
\end{itemize}
All classical predictions are recovered in the continuum limit $N \to \infty$. The framework provides a natural bridge to quantum gravity without invoking additional structures.

\section{Review of Continuous Dual Spacetime}
\label{sec:review}

Each particle carries usual spacetime vectors $X = ct j + x kI + y kJ + z kK$ and dual vectors $X' = ct' k + x' jI + y' jJ + z' jK$, related by $X' = Xi$. The complete rotor is
\[
R_\text{total} = \exp\left( \sum_{a=1}^3 \frac{\omega_a}{2} i\Gamma_a + \frac{\phi_a}{2} \Gamma_a \right), \quad \Gamma_1 = I,\ \Gamma_2 = J,\ \Gamma_3 = K.
\]
The torsional mismatch rotor $\Omega = R_\text{usual}^\dagger R_\text{dual}$ yields the bivector $\Omega_\text{biv} = \log \Omega$, and the scalar
\[
J = \frac12\sum_{a=1}^3(\alpha_a^2-\beta_a^2),
\]
with $\Jnorm=\frac8{3\pi^2}J$. Under the half-angle expansion $\Omega_\text{biv}=\sum_a\bigl(\frac{\alpha_a}{2}i\Gamma_a+\frac{\beta_a}{2}\Gamma_a\bigr)$ and $B(i\Gamma_a,i\Gamma_a)=+8$, $B(\Gamma_a,\Gamma_a)=-8$, one has $B(\Omega_\text{biv},\Omega_\text{biv})=2\sum_a(\alpha_a^2-\beta_a^2)$, so $J$ is four times $\frac1{16}B(\Omega_\text{biv},\Omega_\text{biv})$. The six generators span a Lorentz candidate $\mathfrak{so}(3,1)$, not $\mathfrak{so}(3,1)\oplus\mathfrak{so}(3,1)$.

The same quadratic data determine an unsigned mass
\[
M = \frac12\sum_{a=1}^3(\alpha_a^2+\beta_a^2),\qquad \Mnorm=\frac8{3\pi^2}M.
\]
Vanishing of $J$ does not force vanishing of $M$: the balanced ray $\alpha_a=\beta_a$ has $\Jnorm=0$ with $\Mnorm>0$ whenever the common rapidity is nonzero. Usual--dual interchange sends $J\mapsto -J$ and preserves $M$.

\section{Discretization of Rapidity Parameters}
\label{sec:discretization}

The dual map $X \mapsto Xi$ implies an intrinsic identification of the time direction with a circle: successive applications of the dual map twice return a sign flip ($Xi i = -X$), suggesting compactification. This compactness of the dual identification must not be confused with periodicity of every rotor exponential.

\begin{theorem}[Sectorial periodicity]
\label{thm:periodicity}
Let $B^-_a$ be a cyclic generator ($B^-_a{}^2=-1$) and $B^+_a$ a hyperbolic generator ($B^+_a{}^2=+1$). Then, for every real $\theta$ and every $k\in\mathbb{Z}$,
\[
\exp\bigl((\theta+2\pi k)B^-_a\bigr)=\exp(\theta B^-_a),
\]
while
\[
\exp\bigl((\theta+2\pi)B^+_a\bigr)\neq\exp(\theta B^+_a).
\]
In particular $\exp(2\pi B^-_a)=1$ and $\exp(2\pi B^+_a)\neq 1$.
\end{theorem}

Thus only the elliptic (dual-rotation) sector is intrinsically $2\pi$-periodic. Placing hyperbolic rapidities on a $2\pi$-torus remains a modelling choice, justified by the dual-map compactification rather than by the exponential of a boost.

We nevertheless discretize both sectors uniformly,
\[
\omega_a = \frac{2\pi n_a}{N}, \quad \phi_a = \frac{2\pi m_a}{N}, \quad n_a, m_a \in \mathbb{Z}/N\mathbb{Z},
\]
where $N \gg 1$ is the effective number of states in the compact direction. Physically, $N \propto m c^2 / E_0$ with $E_0$ a fundamental energy (Planckian or particle-specific). Canonical representatives satisfy $0\leq n_a,m_a<N$, hence the embedded rapidities lie in $[0,2\pi)$.

\begin{proposition}
\label{prop:torus-card}
The discrete phase space $(\mathbb{Z}/N\mathbb{Z})^6$ has cardinality $N^6$.
\end{proposition}

The mismatch rotor $\Omega$ therefore takes values in a finite set, namely the image of this torus under the torsion-rotor exponential. Every discrete rotor is unitary: $R\,\widetilde{R}=1$.

\section{Bounded Torsion Scalar on Admissible Configurations}
\label{sec:bounded}

\begin{definition}[Admissible configuration]
\label{def:admissible}
A discrete configuration is \emph{admissible} when the embedded rapidities satisfy $\alpha_a,\beta_a\geq 0$ and $\alpha_a+\beta_a\leq\pi/2$ for each axis $a=1,2,3$ (see Ref.~\cite{diophantine-paper} for the surrounding discussion).
\end{definition}

Admissibility is an extra hypothesis: it is not implied by placing the six rapidities on $(\mathbb{Z}/N\mathbb{Z})^6$. Non-negativity of the embedding is automatic for canonical torus representatives, so the content of the hypothesis is the anti-synchronisation inequality. Equivalently, on each axis,
\begin{equation}
\label{eq:four-le}
4(n_a+m_a)\leq N.
\end{equation}
The principal-branch condition $|\alpha_a+\beta_a|\le\pi/2$ alone does not bound $J$.

The Killing form asymmetry implies that maximal attraction occurs when usual boosts dominate ($\alpha_a \gg \beta_a$), and maximal repulsion when dual rotations dominate. On the continuous admissible cone the extremal values are pure hyperbolic dominance ($\alpha_a=\pi/2$, $\beta_a=0$, all $a$), giving $J=\Jbound$ and $\Jnorm=1$, or pure elliptic dominance ($\alpha_a=0$, $\beta_a=\pi/2$), giving $J=-\Jbound$ and $\Jnorm=-1$. Equality $|\Jnorm|=1$ holds if and only if every axis is of the same extremal type. The unsigned mass obeys the same ceiling: $0\leq\Mnorm\leq 1$, with equality at either wall.

\begin{theorem}[Continuous admissible bound]
\label{thm:continuous-bound}
For every admissible continuous configuration,
\[
|J| \leq \Jbound, \qquad |\Jnorm| \leq 1, \qquad 0\leq \Mnorm\leq 1.
\]
The image of $\Jnorm$ on that cone is exactly $[-1,1]$.
\end{theorem}

Configurations that violate admissibility (e.g.\ negative embedded rapidities while $|\delta_a|\leq\pi/2$) need not satisfy this bound.

\subsection{Integer mismatch and the discrete spectrum}

Write
\[
\Delta(t)=\sum_{a=1}^3\bigl(n_a^2-m_a^2\bigr),\qquad
\Sigma(t)=\sum_{a=1}^3\bigl(n_a^2+m_a^2\bigr)
\]
for a torus point $t=(n_a,m_a)$. The torsional scalars are then exactly the rational multiples
\begin{equation}
\label{eq:spectrum}
\begin{aligned}
J&=\tfrac12\bigl(2\pi/N\bigr)^2\Delta(t),
&
\Jnorm&=\frac{16}{3N^2}\Delta(t),\\
M&=\tfrac12\bigl(2\pi/N\bigr)^2\Sigma(t),
&
\Mnorm&=\frac{16}{3N^2}\Sigma(t).
\end{aligned}
\end{equation}
In particular $\Jnorm=0$ if and only if $\Delta(t)=0$. Balanced axes $n_a=m_a$ give a sufficient condition for vanishing $J$, but not a necessary one: the cross-axis cancellation $(n,m)=(1,0)$, $(0,1)$, $(0,0)$ is admissible for every $N\geq 4$, has $\Delta=0$ and $\Sigma=2$, and is not balanced. Unsigned mass and torsion occupy the even sublattice
\[
|\Delta|\leq\Sigma\leq 3\bigl\lfloor N/4\bigr\rfloor^2,\qquad
\Sigma\equiv\Delta\pmod{2}.
\]
Vacuum is $\Sigma=0$, hence only the origin. Vanishing $J$ with positive mass is $\Delta=0$ and $\Sigma>0$.

The one-axis mismatches are the differences of squares $n^2-m^2$ in the triangle $n+m\leq K=\lfloor N/4\rfloor$. The three-axis image of $\Delta$ is the corresponding three-fold sumset, a subset of $\{-3K^2,\ldots,3K^2\}$. For $N\geq 4$ one has $K\geq 1$, so $\Delta=\pm 1$ is attained (one axis at $(1,0)$ or $(0,1)$, the others at the origin). Consequently the universal lower bound on nonzero height is sharp:
\begin{equation}
\label{eq:gap}
\min\bigl\{|\Jnorm|:\ \text{admissible},\ \Delta\neq 0\bigr\}
=\frac{16}{3N^2}\qquad(N\geq 4).
\end{equation}
The bound interval is not fully occupied. At $N=4$ ($K=1$) every integer in $\{-3,\ldots,3\}$ occurs; already at $N=8$ ($K=2$) the slots $|\Delta|\in\{10,11\}$ are empty, and at $N=12$ ($K=3$) ten slots inside $[-27,27]$ remain vacant (including $|\Delta|\in\{20,23,24,25,26\}$). Density as $N\to\infty$ through multiples of $4$ still fills $[-1,1]$; at finite $N$ the discrete spectrum has holes.

\begin{theorem}[Sharp discrete bound]
\label{thm:sharp}
On every admissible discrete configuration,
\[
\bigl|\Delta(t)\bigr|\leq 3\bigl\lfloor N/4\bigr\rfloor^2,\qquad
\Sigma(t)\leq 3\bigl\lfloor N/4\bigr\rfloor^2,
\]
and therefore
\[
|\Jnorm|\leq\Bigl(\frac{4\lfloor N/4\rfloor}{N}\Bigr)^2,\qquad
\Mnorm\leq\Bigl(\frac{4\lfloor N/4\rfloor}{N}\Bigr)^2\leq 1.
\]
If $4\mid N$, both signs $\Jnorm=\pm 1$ are attained, at the pure-hyperbolic lattice point $n_a=N/4$, $m_a=0$ and the pure-elliptic point $n_a=0$, $m_a=N/4$ respectively. If $4\nmid N$, every admissible point satisfies $|\Jnorm|<1$. Equality in the sharp ceiling holds if and only if all three axes sit on the \emph{same} wall: all $(K,0)$ or all $(0,K)$. Mixed walls cancel and lie strictly inside the bound.
\end{theorem}

\begin{corollary}
\label{cor:small-N}
If $N<4$, the only admissible point is the origin, and $\Jnorm=0$ identically on the admissible cone.
\end{corollary}

For $N\geq 8$ the balanced lattice point $n_a=m_a=1$ (all axes) is admissible, has $\Delta=0$, and has strictly positive mass. The balanced cube itself has cardinality $(\lfloor K/2\rfloor+1)^3$; the zero-height locus is strictly larger once $K\geq 1$ (seven points at $N=4$, thirty-two at $N=8$, sixty-eight at $N=12$). Vanishing $J$ on the discrete cone does not force the vacuum.

\subsection{Cardinality of the admissible cone}

Inequality~\eqref{eq:four-le} factorizes across axes. On each axis the pairs $(n,m)$ with $n+m\leq\lfloor N/4\rfloor$ form a triangular region of cardinality $T(K)=(K+1)(K+2)/2$ with $K=\lfloor N/4\rfloor$.

\begin{theorem}[Admissible cardinality]
\label{thm:card}
The number of admissible configurations is
\[
T\bigl(\lfloor N/4\rfloor\bigr)^3=\Biggl(\frac{\bigl(\lfloor N/4\rfloor+1\bigr)\bigl(\lfloor N/4\rfloor+2\bigr)}{2}\Biggr)^3.
\]
In particular the count depends on $N$ only through $\lfloor N/4\rfloor$. For $N<4$ it equals $1$; for $N=4$ it equals $27$; for $N=8$ it equals $216$.
\end{theorem}

The admissible values of $\Jnorm$ therefore form a finite subset of the scaled integer lattice $\frac{16}{3N^2}\mathbb{Z}$, lying in the interval determined by Theorem~\ref{thm:sharp}, and occupying only those slots whose integer mismatch lies in the three-fold sumset. As $N\to\infty$ through multiples of $4$, this discrete image becomes dense in $[-1,1]$: every value in $[-1,1]$ is approximated arbitrarily closely by $\Jnorm$ of an admissible torus point, even though at each finite $N$ empty slots remain.

In the continuum limit the admissible envelope $\alpha_a,\beta_a\in[0,\pi/2]$ with $\alpha_a+\beta_a\leq\pi/2$ persists; the bound $|\Jnorm|\leq 1$ therefore remains valid in the limiting theory on that domain, while TEGR equivalence is recovered locally.

\section{Integer Biquaternion Ring and Finite Rotor Image}
\label{sec:integer}

Define the integer biquaternion ring as the $\mathbb{Z}$-span of the basis with Gaussian/Hurwitz integer coefficients in each quaternion sector (precise order to be classified). That order is not constructed here, and a Lorentzian integer order may have infinitely many units (Pell-type).

What is finite---and all that the discrete torus supplies---is the \emph{discrete rotor image}: the image of $(\mathbb{Z}/N\mathbb{Z})^6$ under the torsion-rotor map $t\mapsto\exp(\Omega_{\mathrm{torsion}}(t))$. Finiteness of the domain implies finiteness of the image; no enumeration, and no identification with $R^\times$ of an integer order, is required. The image of the admissible subclass is a finite subset of that rotor image.

Same-axis generators commute, so a one-axis mixed exponential factorises as a product of a hyperbolic rotor and an elliptic rotor, with closed forms $\cosh/\sinh$ and $\cos/\sin$. Distinct-axis generators do not commute. A generic six-parameter discrete rotor is therefore not the product of three independent one-axis factors, and the discrete rotor image is not the Cartesian product of three one-axis images.

Quantizing hyperbolic rapidities onto a $2\pi$-torus remains a modelling choice, as Theorem~\ref{thm:periodicity} records. For balanced-type seeds the modular winding number vanishes, so a winding obstruction on that torus cannot fire.

\section{Diophantine Constraints on Gravitational Solutions}
\label{sec:diophantine}

The field equations in the discrete theory reduce to finding integer solutions $(n_a, m_a)$ satisfying Diophantine equations derived from the action principle, searched over the admissible cone of Theorem~\ref{thm:card}. Zero torsional height is the integer condition $\Delta(t)=0$. For example, spherically symmetric solutions correspond to integer points on elliptic curves parameterizing radial torsional profiles.

The layered torsional structure (attractive/repulsive alternations) manifests as finite resonance conditions. Finiteness of the admissible rotor classes at each $N$ is a theorem; a numerical nuclear cutoff is not.

\section{Physical and Observational Implications}
\label{sec:implications}

\begin{itemize}
  \item \textbf{Ultraviolet cutoff}: Maximal $|\Jnorm| = 1$ (equivalently $|J| = 3\pi^2/8$) on admissible configurations implies bounded acceleration, eliminating singularities. At finite $N$ the effective ceiling is the stricter value $(4\lfloor N/4\rfloor/N)^2$, attained only at the same-wall lattice points. The theory also has a floor: the smallest nonzero torsional quantum is $\varepsilon_N=16/(3N^2)$ for every $N\geq 4$.
  \item \textbf{Discrete spectra}: Gravitational wave echoes and nuclear energy levels acquire discrete torsional contributions on the occupied slots of the lattice~\eqref{eq:spectrum}. Not every multiple of $\varepsilon_N$ inside the bound is realised; the empty slots at finite $N$ are part of the spectrum.
  \item \textbf{Nuclear chart}: Finiteness of the discrete rotor image implies only that there are finitely many admissible rotor classes at each $N$, namely at most $T(\lfloor N/4\rfloor)^3$ of them. A numerical cutoff $A\lesssim 300$ is an underived heuristic, not a theorem of the discrete algebra.
  \item \textbf{Experimental tests}: Search for discrete torsional echoes in LIGO/Virgo data; precision clock networks probing planetary cores for exact $J=0$ cancellation. A $J=0$ detection would not by itself imply vanishing mass.
\end{itemize}

\section{Conclusions and Outlook}
\label{sec:conclusions}

The discretization of dual spacetime completes the rejection of the continuum, yielding a bounded torsion scalar on the admissible cone, a sharp discrete ceiling governed by $\lfloor N/4\rfloor$ and attained only at same-wall configurations, an explicit triangular count of admissible points, a universal attained floor $\varepsilon_N=16/(3N^2)$, a three-fold sumset spectrum with holes at finite $N$, a finite discrete rotor image that is not a product of one-axis images, and a Diophantine search space. Cyclic rotor exponentials are $2\pi$-periodic; hyperbolic ones are not, so torus quantization of boosts remains a modelling choice. Future work includes constructing an actual integer order (whose unit group need not be finite) and explicit descent arguments connecting macroscopic GR solutions to Planck-scale integer configurations. Admissibility, the rotor-image finiteness, and any nuclear-chart cutoff remain logically independent of one another.

This framework elevates dual spacetime theory to a candidate for quantum gravity grounded in classical algebraic geometry and number theory.

\bibliographystyle{unsrt}
\begin{thebibliography}{2}
\bibitem{main-paper}
Hypernumbernet Collaboration, ``Gravity as Torsion between Dual Spacetime: A Biquaternionic Reformulation of General Relativity,'' arXiv:xxxx.xxxxx [gr-qc] (2025).
\bibitem{diophantine-paper}
Hypernumbernet Collaboration, ``Discrete Biquaternionic Double Spacetime: Rigorous Proofs of Diophantine Conjectures via Torsional Mismatch and PGA Embedding,'' (2025).
\end{thebibliography}

\end{document}
